% =============================================================================
% Pillar 3 -- Iter 135: Native-Wu test + threshold extrapolation +
% ZVF cross-link. Driver: scripts/group_size_iter135.py.
% Inputs:  experiments/results/groupsize_zvf_sweep.{json,tsv} (n=12 runs,
%          4 G-values x 3 seeds),
%          experiments/results/group_size_token_normalized.tsv (n=20 cells).
% Outputs: experiments/results/group_size_iter135_*.tsv,
%          figures/group_size_iter135.pdf
% =============================================================================
\section{Iter 135 -- Native-Wu test, threshold extrapolation, and the
ZVF retention cross-link}
\label{sec:group-size-iter135}

\paragraph{Setup.}
Iter~131 established that the Wu et al.\ (2025) G$\,{=}\,2 \sim$ G$\,{=}\,16$
claim is \emph{budget-conditional}: on Qwen2.5-0.5B / arithmetic, G$\,{=}\,4$
retention of G$\,{=}\,32$ falls monotonically with $\log_{10}T$ and breaks the
97.6\% Wu bound at every $T\!\ge\!4$M. Iter~127 showed the joint
acc$(G,T)$ surface and the bounded cone at $G^{\ast}\!\le\!32$.
Iter~135 closes three decisive follow-on questions:

  (A) Does Wu's claim hold on its \emph{native} $G\!=\!2$ vs $G\!=\!16$ pairing
      on this benchmark, or does it already fail?
  (B) What is the \emph{extrapolated} compute threshold $T^{\ast}$ at which
      G$\,{=}\,4$ retention of G$\,{=}\,32$ drops below $50\%$?
  (C) How tightly does the Pillar~3 retention collapse track the
      Pillar~2 ZVF mechanism?

\paragraph{Finding 1 -- the Wu et al.\ (2025) claim is supported on its native
$G\!=\!2$ vs $G\!=\!16$ at this benchmark.}
On the $n\!=\!3$ seed-paired runs of the groupsize\_zvf\_sweep, the
heldout-accuracy retention acc$(G\!=\!2)/$acc$(G\!=\!16)$ is
$\hat{R}=1.0035$ with bootstrap 95\% CI $[0.9899, 1.0206]$.
The paired accuracy is acc$(G\!=\!2)=0.982\pm0.008$ vs
acc$(G\!=\!16)=0.978\pm0.010$ (paired diff $-0.003\pm0.009$,
Cohen's $d=-0.22$). G$\,{=}\,2$ is statistically equivalent to (and
marginally above) G$\,{=}\,16$ on this near-ceiling arithmetic task.
\textbf{By the retention and TOST criterion, the Wu claim survives at its
native $G\!=\!2$ vs $G\!=\!16$ scale on this near-ceiling task.}
The 97.6\% retention threshold is comfortably exceeded (TOST@0.976
$p\!<\!10^{-3}$).

\paragraph{Finding 2 -- but the reconstructed grid contradicts the larger $G\!=\!4$ vs
$G\!=\!32$ pairing at every $T\!\ge\!4$M, and the threshold $T^{\ast}$
for the 97.6\% bound is below $T\!=\!10^{6}$ tokens.}
Extrapolating iter~131's OLS retention from the reconstructed token-normalized
grid, $\hat{R}(T) = 1.776 - 0.138 \cdot
\log_{10}T$ (R$\,{=}\,{-}0.952$, $p\!=\!0.048$, $n\!=\!4$):
\begin{itemize}\itemsep2pt
\item $T^{\ast}_{97.6} = 6.27\!\times\!10^{5}$ tokens (Wu-native bound, $<\!1$M)
\item $T^{\ast}_{95} = 9.67\!\times\!10^{5}$ tokens (Wu-strict bound, $\sim\!1$M)
\item $T^{\ast}_{90} = 2.23\!\times\!10^{6}$ tokens (practitioner bound, $\sim\!2$M)
\item $T^{\ast}_{50} = 1.76\!\times\!10^{9}$ tokens (collapse bound, $>\!1$B)
\end{itemize}
The $T^{\ast}_{97.6} \!\approx\! 6.3\!\times\!10^{5}$ prediction is
\emph{just below} the lowest reconstructed grid budget $T\!=\!10^{6}$, consistent
with the iter131 observation that at $T\!=\!1$M retention is
$R\!=\!0.976$ (exactly the Wu bound). Iter~135 sharpens iter131's
``budget-conditional'' framing into a \emph{quantitative threshold}:
the reconstructed grid predicts that the Wu-style retention threshold holds at
$T\!\le\!T^{\ast}_{97.6}$ and fails for any
$T\!\ge\!T^{\ast}_{90}\!\approx\!2.2$M on G$\,{=}\,4$ vs G$\,{=}\,32$.

\paragraph{Finding 3 -- the reward-vs-$G$ curve is FLAT (consistent with
Wu), not steep.}
The 4-point reward curve on the same sweep is best fit by
$\hat{r}(\log_{10}G) = 0.991 - 0.060 \log_{10}G + 0.045 (\log_{10}G)^{2}$
(quadratic, $R^{2}\!=\!0.999$), with linear-only slope
$+0.0072$/decade (95\% bootstrap CI $[-0.0049, +0.0144]$, i.e.\ crosses
zero). The paired reward diff at $G\!=\!16$ vs $G\!=\!2$ is
$+0.0066 \pm 0.0049$, one-sided $p\!=\!0.037$. The reward signal is
weakly monotonically increasing in $G$ with marginal statistical
significance at $n\!=\!3$ paired seeds -- \emph{reward} does not move
much on this near-ceiling task, but the \emph{gradient noise}
reduction matters (the iter123/127 mechanism).

\paragraph{Finding 4 -- the Pillar 3 retention collapse links to, but is not
fully explained by, the Pillar 2 ZVF mechanism.}
The cross-pillar quantification is direct:
\begin{itemize}\itemsep2pt
\item ZVF drops $0.207$ (24.7\%) as $G$ grows $2\!\to\!16$ on the sweep
      (mean ZVF $0.838\!\to\!0.631$, 3 seeds).
\item Retention of G$\,{=}\,4$ over G$\,{=}\,32$ drops $0.249$ as $T$
      grows $1$M$\!\to\!64$M (retention $0.976\!\to\!0.727$, 20 cells).
\item Cross-link ratio $\rho = \Delta R_{T} / \Delta$ZVF$_{G} = 0.249 / 0.207
      = +1.20$.
\end{itemize}
A $\rho\!\approx\!1$ cross-link coefficient means that the same ZVF axis is
associated with both trends, but iter115 shows this is not a complete causal
explanation. The iso-$G$ contrast-yield curve
shows this directly: G$\,{=}\,4$ carries mean\_GU$\,=2.15$ (4.4$\times$ G$\,{=}\,32$'s
$0.49$) but loses $24.7$\% of G$\,{=}\,32$'s accuracy. \textbf{Low $G$
therefore trades high contrast-yield per group against gradient-noise and
amortization costs that the reconstructed grid cannot separate without a
direct matched-budget run.}

\paragraph{Finding 5 -- log-log ZVF slope $\beta = -0.137$/decade is
the canonical empirical signature; the linear $G$-slope is
$-0.230$/decade (consistent with iter131).}
On the $4G\!\times\!3$-seed sweep, $\log_{10}(\text{mean\_zvf}) =
-0.035 - 0.137 \cdot \log_{10}G$ ($R^{2}\!=\!0.9996$), or equivalently
mean\_zvf $= 0.904 - 0.230 \cdot \log_{10}G$ ($R^{2}\!=\!0.997$). The
$|\beta_{\log\log}| = 0.137$ is well below the i.i.d.\ Bernoulli
prediction $|\beta|\!=\!1.0$ because of the anti-herding diversity bonus
$\delta_{\text{div}}$ (iter11/iter131): the empirical ZVF is much
higher than the i.i.d.\ formula
$p^{G} + (1-p)^{G}$ would predict, because high-temperature
autoregressive sampling at fixed prompt anti-herds ($\rho\!<\!0$).

\paragraph{Decisive summary.}
Wu et al.\ (2025) is supported on its native pairing (G$\,{=}\,2$
vs G$\,{=}\,16$). The reconstructed grid predicts that the larger
G$\,{=}\,4$ vs G$\,{=}\,32$ pairing breaks the same retention threshold by
$T\!\ge\!2.2$M tokens, but that extrapolation should be treated as a
hypothesis until a directly measured $G{=}32$ matched-budget cell exists. The
cross-pillar ratio $\rho\!=\!1.20$ ties Pillar~3 to the Pillar~2 ZVF axis,
while iter115's GU-ratio result points to gradient-noise or amortization costs
as the binding mechanism in the reconstructed canonical-scale regime.
